\hypertarget{pd15:investment}{%
\section{The value of obtaining additional evidence}\label{pd15:investment}}

\hypertarget{proposal-acceptance-and-later-qualification}{%
\subsection{The evidence needed before investing}\label{proposal-acceptance-and-later-qualification}}

The investment decision requires a comparison of feasible alternatives,
their costs and the evidence for their expected benefits. Reproducible
public and simulated studies can support that comparison within their
stated assumptions. A restricted result may justify a narrow prototype
while leaving broader research questions unresolved.

Public evidence informs three choices: which existing methods to reuse,
where language assistance is useful, and how analysts should conduct a
reproducible study. A dataset can reveal a timing error or support a
controlled comparison without settling every question about those choices.

Later bank qualification asks a different question: do the assumptions,
precision, workflows and economics hold for a specific institution? That
stage may need genuine offer histories, logging information, independent
labels, owner-reviewed semantics, production constraints and prospective
outcomes. The relevant input should be requested only for the capability
and target population under test. An internal bank pilot is not required
to demonstrate that the proposal has a reproducible public baseline and
a sensible experiment design.

\hypertarget{a-value-of-information-decision}{%
\subsection{A value-of-information
decision}\label{a-value-of-information-decision}}

Let \(d\in\{\text{stop},\text{build},\text{pilot}\}\) denote an
investment choice, \(\theta\) unresolved technical and economic
conditions, and \(U(d,\theta)\) its net value. Let \(E\) denote current
public/prototype evidence and \(Z\) the possible result of a proposed
experiment. A planning expression for the expected value of sample
information is

\[
\operatorname{EVSI}
=\mathbb E_{Z\mid E}\!\left[\max_d\mathbb E\{U(d,\theta)\mid E,Z\}\right]
-\max_d\mathbb E\{U(d,\theta)\mid E\}.
\]

An experiment merits consideration when its expected decision
improvement exceeds its full incremental cost, including data
preparation, customer burden where relevant, security review, labeling,
analyst time and opportunity cost. This is a project decision framework,
not an empirically estimated ROI. Where a probability model for
\(\theta\) is unjustified, use a scenario table and sensitivity bounds
instead of a precise-looking EVSI number.

A simple illustrative calculation shows the sequencing. Suppose, purely
for planning, a small bank study costing 20,000 could avert an
unsuitable 200,000 build in 20\% of scenarios, with no other benefits or
costs. Its gross avoided-spend value would be 40,000 and its net would
be 20,000 under those assumptions. If the study changes the build
decision in only 2\% of scenarios, the same gross value is 4,000 and it
is not justified on that ground. These are stipulated arithmetic
examples, not quotations, bank probabilities or forecasts.

The decision should identify the uncertainty that matters. If a public
benchmark shows repeated action-clock mistakes, improve the formulation
layer before requesting confidential data. If a correct estimator is too
slow on modest data, fix or replace it before commissioning a platform
pilot. If all plausible effects lead to the same action, a costly
precision improvement may have little value. If policy value reverses
when a missing hardship variable is plausible, a carefully scoped bank
measurement study can have high value even when the public-data
implementation is technically sound.

\hypertarget{conditions-for-moving-to-expensive-evidence}{%
\subsection{Conditions for moving to expensive
evidence}\label{conditions-for-moving-to-expensive-evidence}}

An actionable next-experiment recommendation specifies the proposed
intervention, target population, candidate baseline, decision-loss
threshold and the smallest information addition that could alter the
decision. It also states an explicit stop or redesign outcome. A bank
experiment should not be recommended merely to fill a missing-evidence
cell.

For funding, the missing addition may be time-stamped offers and an
assignment mechanism, not more aggregate deposits. For servicing, it may
be feasible-action and hardship information at each decision time. For
discovery, it may be one feasible intervention that distinguishes
competing graphs. For observation construction, it may be a
representative independently resolved sample of uncertain identities
with known inclusion probabilities. Public data help design each
measurement and estimate its computational burden, while the bank
supplies the later population-specific truth.

Authentic institutional approval, independent acceptance, actual
workflow prices and production service levels remain future
qualification requirements. A proposal can be complete without them. It
cannot honestly claim those outcomes without obtaining them.

\hypertarget{remaining-survey-and-prototype-gaps}{%
\subsection{Remaining survey and prototype
gaps}\label{remaining-survey-and-prototype-gaps}}

The public-data portfolio supports specific analyses, including the
earnings experiment and the UCI timing comparison in this volume. These
executions do not qualify every candidate dataset. Current-source web
snapshots still need exact data-file vintages for future analyses.
Several benchmark datasets need data-specific rights checks. Mortgage
histories require registration and use review; the HMDA publication
endpoint was inaccessible in this run. The Criteo license is a material
constraint for a commercial proposal. These affect individual options;
project-generated alternatives and the already local public data keep
the prototype work moving.

No inspected public financial source supplies the complete logged bank
servicing action process or independent real-bank linkage truth. That
finding is limited to the surveyed sources. It calls for transparent
semi-synthetic experiments now and a justified targeted measurement
later, rather than a claim of impossibility. Public reasoning benchmarks
also cannot establish a live agent's quality without an actual
controlled run, nor resolve the business owner's normative choice of
outcome and feasible intervention.

Finally, public data do not repair the open mathematical arguments for
persistent dependence, hidden states, graph-equivalence effect sets or
arbitrary learning procedures. Those questions remain scoped research
obligations. The proposal can present negative findings and restricted
supported regimes, compare alternatives and describe what evidence would
change the recommendation. That is a strong proposal standard without
pretending that an accepted bank product already exists.
